\chapter*{Vorbemerkung}
\addcontentsline{toc}{chapter}{Vorbemerkung}
\label{chap:vorbemerkung}

In dieser Arbeit wird das generische Maskulinum verwendet, wenn im Mittelpunkt der Aussagen keine spezifischen Personen stehen, sondern die Funktion, die Personen ausüben. Sämtliche biologischen Geschlechter und subjektiv empfundenen Zuordnungen sind damit gleichermaßen gemeint.
